Let now $\Omega$ be a bounded domain in $\mathbb{R}^m$ with a smooth boundary.\footnote{Since this section focuses on introducing basic concepts without any precise statement of results, we shall be vague about smoothness assumptions. ``Smooth'' should therefore be interpreted to mean ``as smooth as may be needed.''} Let $L(\bm{x},D)$ be a differential operator of order $2p$ with smooth coefficients defined on $\overline\Omega$. Moreover, let $B_j(\bm{x},D)$, $j=1,\ldots,p$, be differential operators of orders less than $2p$ which are defined for $\bm{x}\in\partial\Omega$. In the following, we are concerned with the boundary-value problem
\begin{equation}\label{eq:general-homogeneous-boundary-value-problem}
\begin{aligned}
L(\bm{x},D)u&=f(\bm{x}),\quad\bm{x}\in\Omega,\\
B_j(\bm{x},D)u&=0,\quad\bm{x}\in\partial\Omega,\quad j=1,\ldots,p.
\end{aligned}
\end{equation}


We assume that there are additional differential operators \[S_j(\bm{x},D),\quad T_j(\bm{x},D),\quad C_j(\bm{x},D),\quad j=1,\ldots,p,\] defined for $\bm{x}\in\partial\Omega$, with the following properties:\begin{enumerate}
\item $S_j$, $T_j$ and $C_j$ have smooth coefficients and orders less than $2p$.
\item Given any set of smooth functions $\phi_j$, $j=1,\ldots,2p$, defined on $\partial\Omega$, there exist functions $u,v\in C^{2p}(\overline\Omega)$ such that on $\partial\Omega$ we have $B_ju=\phi_j$, $S_ju=\phi_{p+j}$, $j=1,\ldots,p$, and, respectively, $C_jv=\phi_j$, $T_jv=\phi_{p+j}$, $j=1,\ldots,p$.
\item For any $u,v\in C^{2p}(\overline\Omega)$, we have
\end{enumerate}
\begin{equation}\label{eq:general-boundary-operator-green-identity}
\begin{aligned}
&\int_\Omega vL(\bm{x},D)u-uL^*(\bm{x},D)v\,\mathrm{d}\bm{x}\\
&\quad=\int_{\partial\Omega}\sum_{j=1}^p S_j(\bm{x},D)uC_j(\bm{x},D)v
-B_j(\bm{x},D)uT_j(\bm{x},D)v\,\mathrm{d}S.
\end{aligned}
\end{equation}


Of course, the question of what assumptions a boundary-value problem must satisfy for such operators to exist is of crucial importance; we shall return to this issue later when we discuss elliptic boundary-value problems. For the moment, we simply take the existence of the $S_j$, $T_j$ and $C_j$ for granted.

\begin{Definition}
\label{def:adjoint-boundary-value-problem}
Let the preceding assumptions hold. Then the boundary-value problem
\begin{equation}\label{eq:general-adjoint-boundary-value-problem}
\begin{aligned}
L^*(\bm{x},D)v&=g(\bm{x}),\quad\bm{x}\in\Omega,\\
C_j(\bm{x},D)v&=0,\quad\bm{x}\in\partial\Omega,\quad j=1,\ldots,p.
\end{aligned}
\end{equation}
is called \textbf{adjoint} to \eqref{eq:general-homogeneous-boundary-value-problem}.
\end{Definition}

We note that if $u$ and $v$ satisfy \eqref{eq:general-homogeneous-boundary-value-problem} and \eqref{eq:general-adjoint-boundary-value-problem}, respectively, then, according to \eqref{eq:general-boundary-operator-green-identity},
\begin{equation}\label{eq:adjoint-solutions-forcing-orthogonality}
\int_\Omega fv-gu\,\mathrm{d}\bm{x}=0.
\end{equation}
We have made no claim that the operators $C_j$ are unique, and indeed, even for ordinary differential equations, the adjoint boundary conditions are determined only up to linear recombination. We can, however, give an intrinsic characterization of the set of functions characterized by the conditions $C_jv=0$.

\begin{Lemma}
\label{lem:adjoint-boundary-intrinsic-characterization}
Let $v\in C^{2p}(\overline\Omega)$ and let $X_B$ denote the set of all $u\in C^{2p}(\overline\Omega)$ such that $B_ju=0$ on $\partial\Omega$ for $j=1,\ldots,p$. Then $(v,Lu)=(L^*v,u)$ for every $u\in X_B$ iff $C_jv=0$ for $j=1,\ldots,p$.
\end{Lemma}

\begin{Proof}
One direction is obvious from \eqref{eq:general-boundary-operator-green-identity}. To see the converse we note that by assumption we can construct $u\in C^{2p}(\overline\Omega)$ such that $B_ju=0$ and $S_ju=\phi_j$, where $\phi_j$ are given smooth functions. If $(v,Lu)=(L^*v,u)$, then \eqref{eq:general-boundary-operator-green-identity} assumes the form
\begin{equation}\label{eq:adjoint-boundary-conditions-test-identity}
\int_{\partial\Omega}\sum_{j=1}^p\phi_j(\bm{x})C_j(\bm{x},D)v\,\mathrm{d}S=0.
\end{equation}
If this holds for arbitrary $\phi_j$, then clearly $C_jv$ must be zero.
\end{Proof}

Thus, although there may be different sets of adjoint boundary conditions, they must be equivalent to each other. As a caution, we note that equivalent sets of boundary conditions need not be linear combinations of each other. For example, let $\partial\Omega$ be a closed curve in $\mathbb{R}^2$ and let $s$ denote arclength. Then the conditions $v=0$ and $\mathrm{d}v/\mathrm{d}s+v=0$ are equivalent, although they are not multiples of each other.

We conclude this subsection with two examples:

\Needspace{4\baselineskip}
\begin{Example}
\label{ex:laplace-dirichlet-neumann-self-adjoint}
We have
\begin{equation}\label{eq:laplace-dirichlet-neumann-green-identity}
\int_\Omega u\Delta v-v\Delta u\,\mathrm{d}\bm{x}=\int_{\partial\Omega}u\dfrac{\partial v}{\partial n}-v\dfrac{\partial u}{\partial n}\,\mathrm{d}S.
\end{equation}
Hence the Dirichlet and Neumann boundary-value problems for Laplace\textquotesingle s equation are their own adjoints.
\end{Example}

\Needspace{4\baselineskip}
\begin{Example}
\label{ex:normal-tangential-adjoint-boundary-problem}
Let $\Omega$ be a bounded domain in $\mathbb{R}^2$ bounded by a closed smooth curve. Let $s$ denote arclength along the curve. Consider the boundary-value problem
\begin{equation}\label{eq:laplace-normal-tangential-boundary-problem}
\Delta u=f(\bm{x}),\ \bm{x}\in\Omega,\qquad\dfrac{\partial u}{\partial n}+\dfrac{\partial u}{\partial s}=0,\ \bm{x}\in\partial\Omega.
\end{equation}
We find
\begin{equation}\label{eq:laplace-normal-tangential-adjoint-calculation}
\begin{aligned}
\int_\Omega u\Delta v-v\Delta u\,\mathrm{d}\bm{x}
&=\int_{\partial\Omega}u\dfrac{\partial v}{\partial n}-v\dfrac{\partial u}{\partial n}\,\mathrm{d}s\\
&=\int_{\partial\Omega}u\left(\dfrac{\partial v}{\partial n}-\dfrac{\partial v}{\partial s}\right)
-v\left(\dfrac{\partial u}{\partial n}+\dfrac{\partial u}{\partial s}\right)\,\mathrm{d}s.
\end{aligned}
\end{equation}
Hence the adjoint boundary-value problem is
\begin{equation}\label{eq:laplace-normal-tangential-adjoint-problem}
\Delta v=g(\bm{x}),\ \bm{x}\in\Omega,\qquad\dfrac{\partial v}{\partial n}-\dfrac{\partial v}{\partial s}=0,\ \bm{x}\in\partial\Omega.
\end{equation}

\end{Example}
